\subsection{Roots on the Calculator} 
When we're working with an irrational root in a concrete application, we often want a \makeblank[1in] approximation of the value. We'll demonstrate this on the TI-84, but any scientific or graphing calculator will have the appropriate functions.
\begin{procedure}{Taking Square Roots on the TI-84}
\begin{enumerate*}
\item Press \TIbut{2nd}\ \TIbutpro{x\textsuperscript{2}} to access the square root function.
\item Enter the number whose root you want to take.
\item Press \TIright\ to move the cursor outside the square root.
\end{enumerate*}
\noi If your calculator has older firmware, or you have a TI-83, you may need to hit \TIbut{)} to end the root, instead.
\end{procedure}
\begin{Example}
Find a decimal approximation of each square root. Round your answer to three decimal places.
\begin{lpic}{}{2}
\foreach \x/\y/\l/\val in {0/1/a/18,1/1/b/94,0/2/c/289,1/2/d/-20} {
	\regtextright{\x*\value{cols}+1}{-\y}{\l. };
	\regtextleft{\x*\value{cols}+1}{-\y}{$\sqrt{\val}\approx{}$};
}
\end{lpic}
\end{Example}
For roots with other indexes, we need to find the appropriate command in the \TI{MATH} menu.
\begin{procedure}{Taking Arbitrary Roots on the TI-84}
\begin{enumerate*}
\item Enter the index.
\item Press \TIbut{math}\ to access the math menu.
\item Use the \TIup\ and \TIdown\ arrows to find \TI{5:\TIsup{x}\textsurd} and press \TIbut{enter}, or press \TIbut{5}.
\item Proceed as for the square root.
\end{enumerate*}
\noi You can also select \TI{4:\TIsup{3}\textsurd(} or press \TIbut{4} for the third root.
\end{procedure}
\begin{Example}
Find a decimal approximation of each root. Round your answer to three decimal places.
\begin{lpic}{}{3}
\foreach \x/\y/\l/\p/\val in {0/1/a/3/18,
															1/1/b/4/21,
															0/2/c/4/81,
															1/2/d/5/-200,
															0/3/e/7/995,
															1/3/f/8/-215} {
	\regtextright{\x*\value{cols}+1}{-\y}{\l. };
	\regtextleft{\x*\value{cols}+1}{-\y}{$\sqrt[\p]{\val}\approx{}$};
}
\end{lpic}
\end{Example}
